% The comb of an arc (Section "Combs" and Lemma "Combs" of v1).
% \input inside a figure; needs \usepackage{tikz} and
% \usetikzlibrary{arrows.meta,decorations.pathreplacing}.
%
% The T-mesh is that of dislocation.tex (b) and separatrix.tex (b): a band cut
% open along the trace tau of the vertex v of valence five, the right side
% glued to the left side two units higher above y = 3 and at the same height
% below; v is drawn at (0,3), (10,3) and (0,5), the vertex c_3 of valence three
% at (1,4), and the empty corner next to c_3 is glued by a quarter-turn about
% c_3 (see dislocation.tex). Every node is at a whole position; check.py
% checks the numbers below.
%
% The arc a is the piece of the separatrix sigma of v from the T-junction
% n_0 = m = (1,5), where the separatrix of c_3 upwards stops on sigma, to the
% T-junction n_1 = q = (10,5), where sigma stops on tau; e_a = (1,0), l_a = 9.
% Foot of n_0: the trace through m perpendicular to a is the separatrix of c_3
% upwards, which starts at c_3: c_0 = c_3. Foot of n_1: the trace through q
% perpendicular to a is tau, which starts at v: c_1 = v. The comb sigma_a runs
% from c_3 up to m (1 unit), along a to q (9 units), and down tau through the
% T-junction n = (10,4) to v (2 units): it passes through T-junctions only.
% Developed along it, p_{c_0} = (1,4) and p_{c_1} = (10,3), and
%   < e_a, p_{c_1} - p_{c_0} > = < (1,0), (9,-1) > = 9 = l_a
% (Lemma "Combs"); the teeth only add the across component -1.
% The feet are distinct, of valences 3 and 5: the comb reads the length of a
% as that of a dislocation (Proposition 4); with two feet of valence 3 it would
% be an exchange (Proposition 8), with a single foot a winding separatrix
% (Proposition 9).
\begin{tikzpicture}[x=0.5cm, y=0.5cm, font=\footnotesize,
    side/.style={draw=black, line width=0.5pt},
    quadgrid/.style={draw=black!18, line width=0.25pt, step={(1,1)}},
    surface/.style={fill=black!3},
    node dot/.style={circle, fill=black, inner sep=0pt, minimum size=2.6pt},
    comb/.style={draw=measured, line width=1.8pt, line join=round},
    arc mark/.style={line width=2.2pt, draw=black},
    dimension/.style={line width=0.5pt, {Bar[width=5pt]}-{Bar[width=5pt]}},
    val3/.pic={\fill[valthree, draw=black, line width=0.3pt]
      (90:3.4pt) -- (210:3.4pt) -- (330:3.4pt) -- cycle;},
    val5/.pic={\fill[valfive, draw=black, line width=0.3pt]
      (90:3.4pt) -- (162:3.4pt) -- (234:3.4pt) -- (306:3.4pt) -- (18:3.4pt) -- cycle;}]
  \definecolor{valthree}{HTML}{E8890C}
  \definecolor{valfive}{HTML}{7C3AAD}
  \definecolor{measured}{HTML}{14934A}

  \begin{scope}
    \fill[surface] (0,1.6) rectangle (10,3) (1,3) rectangle (10,4) (0,4) rectangle (10,8);
    \draw[quadgrid] (0,1.6) grid (10,3);
    \draw[quadgrid] (1,3) grid (10,4);
    \draw[quadgrid] (0,4) grid (10,8);
    % All traces of the T-mesh in the picture.
    \draw[side] (0,1.6) -- (0,3) (10,1.6) -- (10,8) (0,5) -- (0,8);
    \draw[side] (0,3) -- (10,3) (1,4) -- (10,4) (1,4) -- (1,5) (0,5) -- (10,5);
    \draw[side] (1,4) -- (1,3) (1,4) -- (0,4) -- (0,5);
  \end{scope}
  % The comb of a: from its foot c_3 up to m, along a, down tau to its foot v.
  \draw[comb] (1,4) -- (1,5) -- (10,5) -- (10,3);
  \draw[comb, line width=3pt] (1,5) -- (10,5);
  \foreach \p in {(1,3), (0,4), (10,4), (0,6), (1,5), (10,5), (0,7)}
    \node[node dot] at \p {};
  \pic at (1,4) {val3};
  \foreach \p in {(0,3), (10,3), (0,5)} \pic at \p {val5};
  % Names.
  \node[above] at (5.5,5) {$a$};
  \draw[-{Stealth[length=4pt]}, line width=0.5pt] (6.6,5.75) -- (8,5.75)
    node[right, inner sep=1pt] {$e_a$};
  \node[above left=-1pt] at (1,5) {$n_0$};
  \node[above right=-1pt] at (10,5) {$n_1$};
  \node[below right=-1pt] at (1,4) {$c_0 = c_3$};
  \node[right=2pt] at (10,3) {$c_1 = v$};
  \node[left=1pt] at (0,4.95) {$v$};
  \node[left=1pt] at (0,3) {$v$};
  \node[measured, right=1pt] at (10,4.5) {$\sigma_a$};
  \node[right=1pt] at (10,7.4) {$\tau$};
  % The developed feet and the e_a component of their separation.
  \draw[dimension] (1,2.2) -- (10,2.2) node[midway, fill=black!3, inner sep=1.5pt]
    {$\langle e_a,\, p_{c_1} - p_{c_0}\rangle = \ell_a = 9$};
  \draw[black!50, dotted, line width=0.6pt] (1,2.2) -- (1,4) (10,2.2) -- (10,3);
\end{tikzpicture}
